% Generated draft: requires mathematical and source review.
\documentclass[11pt]{article}
\usepackage[margin=25mm]{geometry}
\usepackage{fontspec}
\setmainfont{Latin Modern Roman}
\usepackage{amsmath,amssymb,mathtools,unicode-math}
\setmathfont{Latin Modern Math}
\usepackage{enumitem,needspace,xurl,hyperref}
\setlength{\parindent}{0pt}
\setlength{\parskip}{5pt}
\setlength{\emergencystretch}{3em}
\newcommand{\N}{\mathbb{N}}
\newcommand{\Z}{\mathbb{Z}}
\newcommand{\Q}{\mathbb{Q}}
\newcommand{\R}{\mathbb{R}}
\newcommand{\C}{\mathbb{C}}
\newcommand{\F}{\mathbb{F}}
\newcommand{\A}{\mathbb{A}}
\newcommand{\PP}{\mathbb{P}}
\newcommand{\E}{\mathbb{E}}
\newcommand{\eps}{\varepsilon}
\newcommand{\dd}{\,\mathrm{d}}
\begin{document}
\section*{UnsolvedMath ID: 6700005}
\textbf{Source problem code:} catalogue entry
\subsection{Definitions and mathematical statement}
Let \(Y\) be a closed smooth manifold of dimension \(m\ne4\) which admits no Riemannian metric of everywhere positive scalar curvature. On the band \(X=Y\times[-1,1]\), write \(\partial_-=Y\times\{-1\}\), \(\partial_+=Y\times\{1\}\), and, for a smooth Riemannian metric \(g\) on \(X\), define
\[
 w(g)=\operatorname{dist}_g(\partial_-,\partial_+)
 =\inf\{\operatorname{Length}_g(\gamma):\gamma\text{ joins }\partial_-\text{ to }\partial_+\}.
\]
The source's optimistic question is whether there is a finite constant \(\delta(Y)>0\), independent of \(g\), such that
\[
 \operatorname{Scal}_g\geq n(n-1)\quad\Longrightarrow\quad w(g)<\delta(Y)
\]
for every such metric. No boundary mean curvature condition is stated. The expected conventional normalization is \(n=\dim X=m+1\), but the extracted passage does not define \(n\); this identification needs confirmation.

The source additionally asks for an ideal explicit bound. Its extracted expression is ``\(\delta=2\pi\ \dim(Y)+1\)'', with fraction layout lost. A likely reading is
\[
 \delta(Y)=\frac{2\pi}{m+1},
\]
but this is recorded as an interpretation requiring source verification, not as an unambiguous extracted formula. The source summary separately writes \(2\pi/(n+1)\), without reconciling its use of \(n\) with the scalar curvature normalization. The question retains strict width inequality and the exclusion of dimension four. The intended lower dimension range is not specified.
\subsection{Short English statement}
Does a closed manifold with no positive scalar curvature metric force every sufficiently positively curved product band over it to have uniformly bounded width? Gromov excludes base dimension four and suggests a sharp explicit constant, whose extracted formula needs checking.
\subsection{Sources}
\begin{itemize}
\item Local UnsolvedMath registry, permanent ID 6700005, source code catalogue entry; Misha Gromov, 2017; update 2026-08-06.
\item Gromov, \emph{101 Questions, Problems and Conjectures around Scalar Curvature} (2017), printed tag [?5], occurrence 5, PDF page 2: \url{https://www.ihes.fr/~gromov/wp-content/uploads/2018/08/101-problemsOct1-2017.pdf}. The registry reports access on 2026-07-29, normalized PDF text extraction, and an incomplete public author-hosted PDF. Current status and reuse terms require review.
\item The supplied partial-progress summary mentions Gromov's width questions, the two-sphere lemma, minimal hypersurface barriers, and Schoen--Yau methods, but supplies no exact further theorem or bibliographic reference.
\end{itemize}
\end{document}
